\chapter{Einleitung}
\label{chapter:introduction}

Moderne Fahrzeuge stellen eine Vielzahl von verschiedenen Assistenzsystemen bereit, darunter beispielsweise eine Spurhalteassistenz oder automatisiertes Einparken \cite{detjen2021supporting}. Gleichzeitig verfügen Nutzer*innen nicht immer über ein ausreichendes Verständnis der Potenziale und Grenzen dieser Assistenzsysteme, was zu falscher Nutzung oder einem übermäßigen Vertrauen in unangemessenen Situationen führen kann \cite{detjen2021supporting}. Für eine sichere Nutzung ist daher entscheidend, dass Fahrer*innen ein angemessenes mentales Modell der jeweiligen Assistenzfunktion entwickeln. Feinauer et al. \cite{feinauer2023first} verstehen mentale Modelle als Wissensstrukturen, die u. a. das Verständnis des Systemumfangs, der Funktionalität und der Gründe für das Systemverhalten umfassen. Darüber hinaus heben sie hervor, dass Fahrer*innen für eine sichere Interaktion mit assistierenden oder automatisierten Fahrzeugfunktionen auch eigene Verantwortung sowie die Grenzen des Systems verstehen müssen. Klassische User Manuals erfüllen diese Aufgabe nur bedingt, da sie nicht zwangsläufig vor der ersten Nutzung wirklich gelesen werden oder Inhalte aus Sicht der Fahrer*innen nur schwer in die konkrete Fahrsituation übertragen werden können \cite{feinauer2023first}. Multimodale oder kontextbezogene Ansätze können Abhilfe schaffen, indem Informationen sowohl zeitlich als auch räumlich näher an der eigentlichen Nutzung vermittelt werden \cite{detjen2021supporting}.